\documentclass[11pt,a4paper]{article}
\usepackage[utf8]{inputenc}
\usepackage{graphicx}
\usepackage[left=2.5cm,top=3cm,right=2.5cm,bottom=3cm,bindingoffset=0.5cm]{geometry}
\usepackage{AEDLogica, AEDEspecificacion, AEDTADs}
\usepackage{caratula}
\usepackage{hyperref}
\newlength{\miSangria}
\setlength{\miSangria}{2em}

\titulo{Trabajo práctico 1: Especificacion de Tads}
\subtitulo{Tad AirMiles}

\fecha{\today}

\materia{Algoritmos y Estructuras de Datos}
\grupo{Grupo: Arraydecuatroposiciones} 

\integrante{Javier M. Ramos}{288/20}{ramosjavier276@gmail.com}
\integrante{Nancy Ramos}{1206/24}{Ramosnancydaniela@outlook.com}
\integrante{Elian Carreño}{1420/23}{elianuba2004@gmail.com }
\integrante{Leandro Amodey}{618/24}{leandroamodey@gmail.com}
% \integrante{Apellido, Nombre2}{002/01}{email2@dominio.com} %


% Declaramos donde van a estar las figuras
% No es obligatorio, pero suele ser comodo
\graphicspath{{../static/}} 

% Asi pueden escribir nuevos comandos. 
% Este por ejemplo asegura q los nombres 
% que figuren con una tipografia diferenciada  
\newcommand{\Tipo}[1]{\mathsf{#1}}
% la sintaxis es \newcommand{\nombreDeLaMacro}[cantidadDeParametros]{Lo que va ser remplazado por el macro} 
\newcommand{\norm}[1]{\vert#1\vert}



\begin{document}
\maketitle
\section{Introducción y Desiciones de Diseño}
El presente trabajo practico es una especificacion formal para un sistema de recompensas basado en millas, como ocurre en muchas aerolineas reales.\\
El detalle de los procedimientos se encuentra en la consigna del trabajo, sin embargo agregamos en esta seccion algunas aclaraciones del porque decidimos alguna solucion en lugar de otra. Tambien en la siguiente seccion detallamos los renombres de tipos y estructuras que usamos en el resto del TAD.\\
Por ultimo añadimos comentarios sobre los propios Procedimientos, para los casos donde creimos que valia la pena alguna aclaracion.\\

\subsection{Observadores}
Elegimos utilizar 3 observadores (Usuarios, Historial y Promociones), que consideramos tienen la informacion mas relevante para describir el estado del sistema.\\
En el caso de Usuarios, tiene un ID para distinguir uno de otro, y la categoria, que si bien puede ser calculada a partir del historial, la consideramos suficientemente importante para ser observada.\\
En el caso de Historial, posee el total de transacciones del sistema, y la informacion necesaria para poder derivar otros valores calculables a partir de ellos. (ejemplo: las millas disponibles, las millas totales, quien transfirio mas, etcetera)\\
Por ultimo las Promociones, identificables de manera unica por nombre, refieren a todas las promociones activas, y en el caso de aquellas con tope, el restante hasta que se agoten y sean eliminadas.

\subsection{Agotamiento de Promociones con Tope}
En la acumulacion de millas, detectamos que cuando a un usuario se le aplica una promocion con tope y aun tiene tope disponibles, se le podrian otorgar una cantidad de millas mayor a la disponible debido a la distancia y su factor de categoria. Para ese caso decidimos que se le otorguen al usuario 2 acumulaciones por separado, una utilizando la promocion con las millas de tope disponibles, y el resto utilizando la "NoPromo", es decir recibe un "bonus parcial".

\subsection{Transferencias Reciprocas Infinitas}
En el caso de la transferencia de millas, detectamos que si 2 usuarios se envian y reciben millas mutuamente, el total de millas de los mismos se incrementaria indefinidamente, dandoles la oportunidad de subir de categoria sin realizar vuelos. Si bien esto lo podiamos resolver evitando que las transferencias acumulen totales, al contradecirse con lo que interpretamos de la consigna decidimos darlo como una opcion valida.

\subsection{Categoria Platino}
Los usuarios que llegan a categoria Platino, se indica que no acumulan mas millas, sin embargo no se especifica que ocurre con las millas disponibles, ni con las transferencias. Nuestra interpretacion de esto, es que estos usuarios SI pueden hacer esas operaciones, siendo unicamente la acumulacion por vuelo la que se anula, utilizando un factor 0 en el calculo de millas otorgadas.

\section{Renombres de Tipos y Estructuras}

\noindent enum Categoría \textbf{ES} \{Bronce, Plata,Oro, Platino  \}

\noindent Transferencia \textbf{ES} Struct $\tupla{
		UsuarioOrigen: \Tipo{\Z},
		UsuarioDestino: \Tipo{\Z},
		Millas: \Tipo{\R}
	}$

\noindent Vuelo \textbf{ES} Struct $\tupla{
		Usuario: \Tipo{\Z},
		Millas: \Tipo{\R},
		Distancia: \Tipo{\Z},
		NombrePromo: \Tipo{\str}, \\
		FactorPromo: \Tipo{\R}
	}$

\noindent Canjeo  \textbf{ES} Struct $\tupla{
		Usuario: \Tipo{\Z},
		Millas: \Tipo{\R}
	}$

\noindent Referido  \textbf{ES} Struct $\tupla{
		UsuarioReferidor: \Tipo{\Z},
		UsuarioReferido: \Tipo{\Z},
		Millas: \Tipo{\R}
	}$

\vspace{2em}

\section{Especificación del TAD AirMiles}
\begin{tad}{Airmiles}


	\comentario{Decidimos dejar en usuarios la categoria por mas que se pueda calcular por fuera. Las promociones tienen a un diccionario que adentro tiene un struct }
	\\
	\obs{Usuarios}{\dict{id:\Tipo{\Z}}{CategoriaActual : Categoria}}
	\\
	\noindent \obs{Vuelos}{\seq{Vuelo}} \\
    \noindent \obs{Transferencias}{\seq{Transferencia}} \\
    \noindent \obs{Canjeos}{\seq{Canjeo}} \\
    \noindent \obs{Referidos}{\seq{Referido}} \\
    \noindent \obs{PromocionesConTope}{\dict{NombrePromocion: \Tipo{\str}}{\\ \struct{FactorPromocion: \Tipo{\R},TopePromocion: \Tipo{\R},MillasRestantes: \Tipo{\R}}}} \\
    \noindent \obs{PromocionesSinTope}{\dict{NombrePromocion: \Tipo{\str}}{FactorPromocion: \Tipo{\R}}} \\

	\phantomsection\hypertarget{sec:crearSistema}{}\label{sec:crearSistema}
	\begin{proc}{crearSistema}
		{}
		{AirMiles}
		\requiere{\Tipo{True}}
		\aseguraLargo{ \norm{res.Usuarios} = 0 \land
			\\
            \norm{res.Vuelos} = 0 \land
            \\
            \norm{res.Transferencias} = 0 \land
            \\
            \norm{res.Canjeos} = 0 \land
            \\
            \norm{res.Referidos} = 0 \land
            \\
            \norm{res.PromocionesConTope} = 0 \land
            \\
			res.PromocionesSinTope = setKey(\dict{}{},"NoPromo", 1)}
	\end{proc}

	\phantomsection\hypertarget{sec:registrarUsuario}{}\label{sec:registrarUsuario}
	\begin{proc}{registrarUsuario}
		{
			\Inout AM: Airmiles,
			\In id: \Tipo{\Z}
		}
		{}
		\comentario{Actualiza el diccionario de usuarios con  nuevo usuario} \\
		\requiere{AM = AM_0 \land id \notin claves(AM.Usuarios) \land id > 0}
		\aseguraLargo{ AM.Usuarios = SetKey(AM_0.Usuarios, id, Bronce) \land \\
			AM.Vuelos = AM_0.Vuelos \land \\ 
            AM.Transferencias = AM_0.Transferencias \land \\
            AM.Canjeos = AM_0.Canjeos \land \\
            AM.Referidos = AM_0.Referidos \land \\
			AM.PromocionesConTope = AM_0.PromocionesConTope \land \\ 
            AM.PromocionesSinTope = AM_0.PromocionesSinTope
		}
	\end{proc}


	\phantomsection\hypertarget{sec:registrarUsuarioConReferidor}{}\label{sec:registrarUsuarioConReferidor}
	\begin{proc}{registrarUsuarioConReferidor}
		{
			\Inout AM: Airmiles,\\
			\In idReferidor: \Tipo{\Z},
			\In idUsuario: \Tipo{\Z}
		}
		{}
		\requiere{AM_0 = AM \land idReferidor \in claves(AM.Usuarios) \land idUsuario \notin claves(AM.Usuarios) \land idUsuario > 0}
		\aseguraLargo{
			% Creo el usuario
			AM.Usuarios = SetKey(AM_0.Usuarios,idUsuario, Bronce) \land
			\\
			\hyperref[sec:EsNuevoReferido]{EsNuevoReferido}(AM,AM_0,idReferidor, idUsuario, 500) \yLuego
			\\
            \comentario{Se chequea si despues de referir a alguien subio de Categoria} \\
            AM.Usuarios[idReferidor] = \hyperref[sec:DarCategoria]{DarCategoria}(AM,idReferidor) \land \\
			AM.PromocionesConTope = AM_0.PromocionesConTope \land \\ 
            AM.PromocionesSinTope = AM_0.PromocionesSinTope \land \\
            AM.Vuelos = AM_0.Vuelos \land \\
            AM.Transferencias = AM_0.Transferencias \land \\
            AM.Canjeos = AM_0.Canjeos \land \\ 
		}
	\end{proc}
	% ===============Aca comienza la funcion de acumular millas===============

	\vspace{1cm}
	\noindent \comentario{Acumula millas al usuario si la categoria del mismo no es platino, en caso de que el nombre de la promo no alcance para completar el viaje del usuario, se crean dos transacciones, una con las millas restantes de la promo y otra con "NoPromo" para completar el viaje. }
	\phantomsection\hypertarget{sec:AcumularMillas}{}\label{sec:AcumularMillas}
	\\
	\begin{proc}{AcumularMillas}
		{
			\Inout AM: \Tipo{Airmiles},
			\In idUsuario: \Tipo{\Z},
			\In distancia: \Tipo{\R},
			\In nombrePromo: \Tipo{\str},
		}
		{}

		\requiere{AM = AM_0 \land idUsuario \in claves(AM.Usuarios) \land nombrePromo \in claves(AM.Promociones) \land distancia > 0 }
		\aseguraLargo{
			\comentario{Si el tope de las millas de la promo supera la distancia se hace una sola transaccion}
			\\
			AM_0.Promociones\lbrack nombrePromo \rbrack.millasRestantes \geq
			\\
			\hyperref[sec:MillasCalculadas]{MillasCalculadas}(distancia, \hyperref[sec:DarFactorCategoria]{DarFactorCategoria}(AM_0.Historial,idUsuario),
			\\ AM_0.Promociones[nombrePromo].factorDePromocion)
			\lor
			\\
			\\
			\comentario{Si la promo no tiene tope lo manejamos a parte}
			\\
			AM_0.Promociones\lbrack nombrePromo \rbrack.millasRestantes = -1
			\implicaLuego
			\\
			\comentario{Se crea la transaccion}
			\\
			\hyperref[sec:EsNuevaTransaccion]{EsNuevaTransaccion}(AM,AM_0,idUsuario,-1,0,
			\\
			\hyperref[sec:MillasCalculadas]{MillasCalculadas}(distancia, \hyperref[sec:DarFactorCategoria]{DarFactorCategoria}(AM_0.Historial,idUsuario),
			\\ AM_0.Promociones[nombrePromo].factorDePromocion) ,distancia, \\
			nombrePromo, AM_0.Promociones[nombrePromo].FactorDePromocion) \land\\
			\\
			\comentario{Se restan las millas de la promocion unicamente para promos que no tiene tope}
			\\
			AM_0.Promociones\lbrack nombrePromo \rbrack.millasRestantes > 0 \implicaLuego
			\\
			AM.Promociones = setKey(AM_0.Promociones, nombrePromo,\\
			AM_0.Promociones[nombrePromo].millasRestantes - \hyperref[sec:MillasCalculadas]{MillasCalculadas}\\(distancia, \hyperref[sec:DarFactorCategoria]{DarFactorCategoria}(AM_0.Historial, idUsuario), \\
			AM_0.Promociones[nombrePromo].factorDePromocion)
			\lor
			\\
			\\
			AM_0.Promociones\lbrack nombrePromo \rbrack.millasRestantes <
			\\
			\hyperref[sec:MillasCalculadas]{MillasCalculadas}(distancia, \hyperref[sec:DarFactorCategoria]{DarFactorCategoria}(AM_0.Historial, idUsuario),\\
			AM_0.Promociones[nombrePromo].factorDePromocion) \implica
			\\
			\\
			\comentario{Se hace 1 transaccion con las millas restantes de promo}
			\\
			\hyperref[sec:EsNuevaTransaccion]{EsNuevaTransaccion}(AM,AM_0,idUsuario,-1,1,
			\\
			AM_0.Promociones[nombrePromo].millasRestantes, distancia,
			\\ nombrePromo, AM_0.Promociones[nombrePromo].FactorDePromocion) \yLuego
			\\
			\\
			\comentario{Se hace otra transaccion con noPromo para completar}
			\\
			\hyperref[sec:EsNuevaTransaccion]{EsNuevaTransaccion}(AM,AM_0,idUsuario,-1,1,
			\\
			MillasACompletar(AM_0.Promociones,"NoPromo",distancia,
			\\\hyperref[sec:DarFactorCategoria]{DarFactorCategoria}(AM_0.Historial,idUsuario), 1)
			\\distancia, "NoPromo", 1) \yLuego
			% Luego de hacer la transaccion tengo que verificar si se modifico la categoría del usuario
			\\
			\comentario{Se puede indefinir si no le pasamos un usuario que este en el sistema}
			\\
			\hyperref[sec:SubioCategoria]{SubioCategoria}(idUsuario,AM.Historial,AM_0.Historial) \implicaLuego
			\\
			AM.Usuarios = SetKey(AM_0.Usuarios,idUsuario,
			\\
			\hyperref[sec:DarCategoria]{DarCategoria}(AM_0.Historial, idUsuario)) \land
			\\
			\comentario{Despues de completar el viaje se verifica si la promocion termino}
			\\
			AM_0.Promociones[nombrePromo].millasRestantes = 0 \implicaLuego
			\\
			\hyperref[sec:SeConsumioLaPromocion]{SeConsumioLaPromocion}(nombrePromo, AM, AM_0)
		}
	\end{proc}

	\vspace{1cm}
	\phantomsection\hypertarget{sec:CanjearMillas}{}\label{sec:CanjearMillas}
	\begin{proc}{CanjearMillas}
		{\Inout AM: AirMiles, \In MillasCanjeadas : \Tipo{\R}, \In IdUsuario : \Tipo{\Z}}
		{}
		\requiereLargo{ AM=AM_0 \yLuego MillasCanjeadas > 0 \land IdUsuario \in claves(AM.Usuarios)}
		\aseguraLargo{
			\IfThenElse{ \hyperref[sec:MillasDisponibles]{MillasDisponibles}(IdUsuario, AM_0.Historial) \geq MillasCanjeadas}
			{\\ \hyperref[sec:EsNuevoCanjeo]{EsNuevoCanjeo}(AM,AM_0,IdUsuario,MillasCanjeadas)}
			{\\AM.Canjeos =  AM_0.Canjeos} \land \\ 
			AM.Usuarios = AM_0.Usuarios \land \\ 
			AM.Vuelos = AM_0.Vuelos \land \\ 
            AM.Transferencias = AM_0.Transferencias \land \\ 
            AM.Canjeos = AM_0.Canjeos \land \\ 
            AM.Referidos = AM_0.Referidos \land \\
            AM.PromocionesConTope = AM_0.PromocionesConTope \land \\
            AM.PromocionesSinTope = AM_0.PromocionesSinTope
		}
	\end{proc}

	\vspace{1cm}
    \comentarioLargo{Aca al hacer la transferencia de millas al final entonces chequeamos la categoría o no? Porque en teoria las transferencias habiamos dicho que no modificaban las totales aunque eso no rompia la lógica, arreglamos/resolvemos eso y no chequeamos el ascenso de categoria o lo hacemos igualmente?}
	\phantomsection\hypertarget{sec:TransferirMillas}{}\label{sec:TransferirMillas}
	\begin{proc}{TransferirMillas}
		{
			\Inout AM: Airmiles,\\
			\In idUsuarioPagador: \Tipo{\Z},
			\In idUsuarioBeneficiario: \Tipo{\Z},
			\In CantMillasTrans: \Tipo{\R}
		}
		{}
		\requiereLargo{AM_0 = AM \land idUsuarioPagador \in claves(AM.usuarios)\\
			idUsuarioBeneficiario \in claves(AM.usuarios) \land \\
			idUsuarioPagador \neq idUsuarioBeneficiario \land \\
			\hyperref[sec:MillasDisponibles]{MillasDisponibles}(AM, idUsuarioPagador) \geq CantMillasTrans
		}
		\aseguraLargo{\hyperref[sec:EsNuevaTransferencia]{EsNuevaTransferencia}( AM, AM_0, idUsuarioPagador, idUsuarioBeneficiario) \land 
        AM.Usuarios = AM_0.Usuarios \land \\ 
        AM.Vuelos = AM_0.Vuelos \land \\ 
        AM.Canjeos = AM_0.Canjeos \land \\ 
        AM.Referidos = AM_0.Referidos \land \\
        AM.PromocionesConTope = AM_0.PromocionesConTope \land \\
        AM.PromocionesSinTope = AM_0.PromocionesSinTope
		}

	\end{proc}

	\vspace{2cm}
	\phantomsection\hypertarget{sec:ConsultarSaldo}{}\label{sec:ConsultarSaldo}
	\begin{proc}{ConsultarSaldo}
		{\In AM: AirMiles, \In IdUsuario : \Tipo{\Z}}
		{\Tipo{\Z}}
		\requiere{IdUsuario \in claves(AM.Usuarios) }

		\asegura{
			res = \hyperref[sec:MillasDisponibles]{MillasDisponibles}(IdUsuario,AM.Historial)
		}
	\end{proc}

	\vspace{1cm}
	\phantomsection\hypertarget{sec:ConsultarCategoria}{}\label{sec:ConsultarCategoria}
	\begin{proc}{ConsultarCategoria}
		{
			\In AM: Airmiles,
			\In idUsuario: \Tipo{\Z},
		}
		{\Tipo{categoria}}
		\requiere{idUsuario \in claves(AM.usuarios)}
		\asegura{res = AM.Usuarios[idUsuario]}

	\end{proc}

	\vspace{1cm}
	\phantomsection\hypertarget{sec:RankingMillas}{}\label{sec:RankingMillas}
	\begin{proc}{RankingMillas}
		{\In AM: Airmiles, \In n : \Tipo{\Z}}
		{\Tipo{\seq{\Tipo{\Z}}}}
		\requiere{ n > 0}

		\aseguraLargo{ (|AM.Usuarios| > n  \land |res| = n) \\ \lor (|AM.Usuarios| \leq n  \land |res| = |AM.Usuarios|) \yLuego
			\\ \hyperref[sec:perteneceUsuarios]{perteneceUsuarios}(AM.Usuarios, res) \land
			\\ \hyperref[sec:sinRepetidos]{sinRepetidos}(res) \land
			\\ \hyperref[sec:ordenDecreciente]{ordenDecreciente}(AM.Historial, res) \land
			\\ \hyperref[sec:losMejoresCandidatos]{losMejoresCandidatos}(AM.Usuarios, AM.
			Historial, res)
		}
	\end{proc}

	

	\vspace{1cm}
	\phantomsection\hypertarget{sec:IniciarPromocion}{}\label{sec:IniciarPromocion}
	\begin{proc}{IniciarPromocion}
		{
			\Inout AM: \Tipo{Airmiles},
			\In NombrePromo: \Tipo{\seq{char}},
			\In FactorPromo: \Tipo{\R},
		}
		{}
		\requiereLargo{
			AM = AM_0 \land
			factorPromo \geq 1 \land \\
			nombrePromo \notin claves(AM_0.PromocionesSinTope)}
		\aseguraLargo{
            AM.PromocionesSinTope = SetKey(AM_0.PromocionesSinTope,\\ nombrePromo, FactorPromo) \land \\
            AM.Usuarios = AM_0.Usuarios \land \\
            AM.Transferencias = AM_0.Usuarios \land \\ 
            AM.Vuelos = AM_0.Vuelos \land \\ 
            AM.Canjeos = AM_0.Canjeos \land \\ 
            AM.Referidos = AM_0.Referidos \land \\
            AM.PromocionesConTope = AM_0.PromocionesConTope			
		}
	\end{proc}
	\vspace{1.0cm}

	\phantomsection\hypertarget{sec:IniciarPromocionConTope}{}\label{sec:IniciarPromocionConTope}
	\begin{proc}{IniciarPromocionConTope}
		{
			\Inout AM: \Tipo{Airmiles},
			\In nombrePromo: \Tipo{\str}, \\
			\In factorPromo: \Tipo{\R},
			\In TopePromo: \Tipo{\R}
		}
		{}
		\requiereLargo{
			AM = AM_0 \land factorPromo \geq 1 \land
            \\ 
            nombrePromo \notin
			claves(AM.PromocionesConTope) \land TopePromo > 0
		}
		\aseguraLargo{
			AM.PromocionesConTope = SetKey(AM_0.PromocionesConTope, nombrePromo,
			\\
			\struct{factorPromo,topePromo,topePromo}) \land
            AM.Usuarios = AM_0.Usuarios \land \\
            AM.Vuelos = AM_0.Vuelos \land \\ 
            AM.Transferencias = AM_0.Usuarios \land \\  
            AM.Canjeos = AM_0.Canjeos \land \\ 
            AM.Referidos = AM_0.Referidos \land \\
            AM.PromocionesSinTope = AM_0.PromocionesSinTope	
		}
	\end{proc}
	% Asegura{ AM.Usuarios = AM_0.Usuarios  ∧ 
	% 		   AM.Historial = AM_0.Historial ∧
	% 		   AM.Promociones = SetKey(AM_0.Promociones, nombrePromo, [factorPromo,0,0] 		
	% 		 }

	\vspace{1.0cm}

	\phantomsection\hypertarget{sec:FinalizarPromocion}{}\label{sec:FinalizarPromocion}
	\begin{proc}{FinalizarPromocion}
		{
			\Inout AM: \Tipo{Airmiles},
			\In nombrePromo: \Tipo{\seq{char}}
		}
		{}
		\requiereLargo{
			AM = AM_0 \land nombrePromo \notin AM_0.PromocionesSinTope \yLuego
			\\
			nombrePromo \in claves(AM.PromocionesConTope)
		}
		\aseguraLargo{
			AM.PromocionesConTope = DelKey(AM_0.PromocionesConTope, nombrePromo) \land
			AM.Usuarios = AM_0.Usuarios \land \\
            AM.Vuelos = AM_0.Vuelos \land \\ 
            AM.Transferencias = AM_0.Usuarios \land \\  
            AM.Canjeos = AM_0.Canjeos \land \\ 
            AM.Referidos = AM_0.Referidos \land \\
            AM.PromocionesSinTope = AM_0.PromocionesSinTope
		}
	\end{proc}

	\vspace{1.0cm}

	\phantomsection\hypertarget{sec:ElQueTieneMasMillas}{}\label{sec:ElQueTieneMasMillas}
	\begin{proc}{ElQueTieneMasMillas}
		{
			\In AM: \Tipo{Airmiles}
		}
		{\Tipo{\Z}}
		\requiere{ \norm{AM.Usuarios} > 0}
		\aseguraLargo{
			\existe{i}{\Tipo{\Z}}{i \in claves(AM.Usuarios) \yLuego res = i } \land
			\\
			\paraTodo{j}{\Tipo{\Z}}{j \in claves(AM.Usuarios) \implicaLuego \hyperref[sec:MillasTotales]{MillasTotales}(AM,res) \geq \hyperref[sec:MillasTotales]{MillasTotales}(AM,j)}
		}
	\end{proc}

	\vspace{1.0cm}



	\phantomsection\hypertarget{sec:ElMasCanjeador}{}\label{sec:ElMasCanjeador}
	\begin{proc}{ElMasCanjeador}
		{
			\In AM: \Tipo{Airmiles}
		}
		{\Tipo{\Z}}
		\requiere{
			\norm{AM.Usuarios} > 0 \land \norm{AM.Historial} > 0
		}
		\aseguraLargo{
			\existe{id}{\Tipo{\Z}}{id \in claves(AM.Usuarios) \yLuego res = id
				\\
				\land \paraTodo{id2}{\Tipo{\Z}}{id2 \in claves(AM.Usuarios) \implicaLuego
					\\
					\hyperref[sec:MillasPerdidas]{MillasPerdidas}(AM,res) \ge
					\hyperref[sec:MillasPerdidas]{MillasPerdidas}(AM,id2)}}
		}
	\end{proc}

	\vspace{1.cm}
	\phantomsection\hypertarget{sec:paresTransferidoresExclusivos}{}\label{sec:paresTransferidoresExclusivos}
	\begin{proc}{paresTransferidoresExclusivos}
		{\In AM: Airmiles}
		{\seq{\tupla{\Tipo{\Z}, \Tipo{\Z}}}}
		\requiere{True}

		\aseguraLargo{
			\hyperref[sec:EstanTodosLosPares]{EstanTodosLosPares}(AM.usuarios,res)
			\\
			\hyperref[sec:sinRepetidos]{sinRepetidos}(res) \land_L
			\\
			\hyperref[sec:NoHayPermutacion]{NoHayPermutacion}(res) \land_L
			\\
			\hyperref[sec:TieneReciprocoHistorial]{TieneReciprocoHistorial}(AM.Historial, res) \land_L
			\\
			\hyperref[sec:SonParejaExclusiva]{SonParejaExclusiva}(AM.Historial, res)
		}
	\end{proc}

    \vspace{0.5cm}
	\phantomsection\hypertarget{sec:ResumenTransacciones}{}\label{sec:ResumenTransacciones}
	\begin{proc}{ResumenTransacciones}
		{\In AM: Airmiles, \In id : \Tipo{\Z}}
		{\tupla{\Tipo{\R}, \Tipo{\R}, \Tipo{\R}}}
		\requiere{ id \in claves(AM.Usuarios)}

		\aseguraLargo{ res_{0} = \hyperref[sec:SumaVuelos]{SumaVuelos}(AM.Vuelos, id) \land
			\\
			res_{1} = \hyperref[sec:SumaCanjeos]{SumaCanjeos}(AM.Canjeos, id) \land
			\\
			res_{2} = \hyperref[sec:SumaTransferenciasRecibidas]{SumaTransferenciasRecibidas}(AM.Transferencias, id) - \\
            \hyperref[sec:SumaTransferenciasEnviadas]{SumaTransferenciasEnviadas}(AM.Transferencias, id) 
		}
	\end{proc}

	\vspace{1cm}
	\phantomsection\hypertarget{sec:MillasTransferidasYRecibidas}{}\label{sec:MillasTransferidasYRecibidas}
	\auxLargo{MillasTransferidasYRecibidas}{Historial:\Tipo{\seq{Transaccion}}, id: \Tipo{\Z}}{\Tipo{\R}}
	{
		\sum\limits_{0}^{\norm{Historial}-1}\IfThenElse{EsTransferenciaEnviada}{-Historial[i].Millas}
		\\{\IfThenElse{EsTransferenciaRecibida}{Historial[i].Millas}{0}}
	}
	\vspace{1cm}
	\predLargo{EsTransferenciaRecibida}
	{Transaccion: \Tipo{Transaccion}, idDestino: \Tipo{\Z}, Usuarios: \Tipo{\dict{\Z}{Categoria}}}
	{
		Transaccion.UsuarioOrigen \in claves(Usuarios) \land
		\\
		Transaccion.TipoDeTransaccion = 1 \land
		\\
		Transaccion.UsuarioDestino = idDestino
	}

    \vspace{0.5cm}
    \phantomsection\hypertarget{sec:EsNuevoReferido}{}\label{sec:EsNuevoReferido}
    \predLargo{EsNuevoReferido}
	{AM: \Tipo{Airmiles}, AM_0: \Tipo{Airmiles}, idReferidor: \Tipo{\Z}, idReferido: \Tipo{\Z}, Millas: \Tipo{\Z}}
	{
		AM.Referidos = AM_0.Referidos ++ [idReferidor,idReferido,Millas]
	}

    \vspace{0.5cm}
    \phantomsection\hypertarget{sec:EsNuevoVuelo}{}\label{sec:EsNuevoVuelo}
    \predLargo{EsNuevoVuelo}
	{AM: \Tipo{Airmiles}, AM_0: \Tipo{Airmiles}, IdUsuario: \Tipo{\Z}, Millas: \Tipo{\R}, distancia: \Tipo{\Z}, NombrePromo: \Tipo{\str}, FactorPromo: \Tipo{\Z}}
	{
		AM.Vuelos = AM_0.Vuelos ++ [IdUsuario,Millas,distancia,NombrePromo,FactorPromo]
	}

    \vspace{0.5cm}
    \phantomsection\hypertarget{sec:EsNuevoCanjeo}{}\label{sec:EsNuevoCanjeo}
    \predLargo{EsNuevoCanjeo}
	{AM: \Tipo{Airmiles}, AM_0: \Tipo{Airmiles}, IdUsuario: \Tipo{\Z}, Millas: \Tipo{\R}}
	{
		AM.Canjeo = AM_0.Canjeo ++ [IdUsuario,Millas]
	}

    \vspace{0.5cm}
    \phantomsection\hypertarget{sec:EsNuevaTransferencia}{}\label{sec:EsNuevaTransferencia}
    \predLargo{EsNuevaTransferencia}
	{AM: \Tipo{Airmiles}, AM_0: \Tipo{Airmiles}, IdUsuarioOrigen: \Tipo{\Z} \\, IdUsuarioDestino: \Tipo{\Z}, Millas: \Tipo{\R}}
	{
		AM.Transferencias = AM_0.Transferencias ++ [IdUsuarioOrigen, \\IdUsuarioDestino,Millas]
	}

	\predLargo{EsTransferenciaEnviada}
	{Transaccion: \Tipo{Transacciones}, idOrigen: \Tipo{\Z}, Usuarios: \Tipo{\dict{\Z}{Categoria}}}
	{
		Transaccion.UsuarioOrigen = idOrigen \land
		\\
		Transaccion.TipoDeTransaccion = 1 \land
		\\
		Transaccion.UsuarioDestino \in claves(Usuarios)
	}

	\vspace{1cm}
	\phantomsection\hypertarget{sec:perteneceUsuarios}{}\label{sec:perteneceUsuarios}
	\predLargo{perteneceUsuarios}{usuarios:\Diccionario{\Tipo{\Z}}{categoria}, res: \seq{\Z}} {
		\paraTodo{i}{\Tipo{\Z}}{0 \leq i < \norm{res} \implica_L res[i] \in claves(usuarios)}
	}

	\vspace{1cm}
	\phantomsection\hypertarget{sec:sinRepetidos}{}\label{sec:sinRepetidos}
	\predLargo{sinRepetidos}{res: \seq{\Z}} {
		\paraTodo{x}{\Tipo{\Z}}{x \in res\implica_L \hyperref[sec:cantAparicion]{cantAparicion}(res,x) = 1}
	}

	\vspace{1cm}
	\phantomsection\hypertarget{sec:cantAparicion}{}\label{sec:cantAparicion}
	\auxLargo{cantAparicion}
	{res: \Tipo{\seq{\Z}}, x: \Tipo{\Z}}
	{\Tipo{\R}}
	{
		\sum\limits_{i=0}^{\norm{res} - 1}
		\IfThenElse{
			res[i] = x
		}{
			1
		}{
			0
		}
	}

	\vspace{1cm}
	\phantomsection\hypertarget{sec:ordenDecreciente}{}\label{sec:ordenDecreciente}
	\predLargo{ordenDecreciente}{Historial:\seq{Transacciones}, res: \seq{\Z}} {
		\paraTodo{i}{\Tipo{\Z}}{0 \leq i < \norm{res}-1 \implica_L \\ millasDisponibles(res[i], Historial) >= millasDisponibles(res[i + 1], Historial)}
	}
	\\
	\phantomsection\hypertarget{sec:losMejoresCandidatos}{}\label{sec:losMejoresCandidatos}
	\\
	\predLargo{losMejoresCandidatos}{usuarios:\Diccionario{\Tipo{\Z}}{categoria}, \\ Historial:\seq{Transacciones}, res: \seq{\Z}} {
		\paraTodo{i}{\Tipo{\Z}}{0 \leq i < \norm{res} \land  res[i] \in claves(usuarios)  \implica_L  \\ \hyperref[sec:elMejorCandidato]{elMejorCandidato}(Usuarios, Historial, res, i) }
	}
	\\
	\phantomsection\hypertarget{sec:elMejorCandidato}{}\label{sec:elMejorCandidato}
	\\
	\predLargo{elMejorCandidato}{usuarios:\Diccionario{\Tipo{\Z}}{categoria}, \\ Historial:\seq{Transacciones}, res: \seq{\Z}, i: \Tipo{\Z}} {
		\paraTodo{x}{\Tipo{\Z}}{ (x \in claves(usuarios) \land \\ millasDisponibles(x,Historial) \leq millasDisponibles(res[i],Historial))\implica_L  \\ (\existe {j}{\Tipo{\Z}}{0 \leq j < i } \land_L res[j] = x) }
	}

	\vspace{1cm}
	\phantomsection\hypertarget{sec:SubioCategoria}{}\label{sec:SubioCategoria}
	\predLargo
	{SubioCategoria}
	{idUsuario: \Tipo{\Z},Historial: \Tipo{Transaccion},Historial0: \Tipo{Transaccion}}
	{\hyperref[sec:DarCategoria]{DarCategoria}(Historial0,idUsuario) \neq \hyperref[sec:DarCategoria]{DarCategoria}(Historial,idUsuario)}


	\vspace{1cm}
	\phantomsection\hypertarget{sec:MillasCalculadas}{}\label{sec:MillasCalculadas}
	\auxLargo
	{MillasCalculadas}
	{distancia:\Tipo{\Z}, Factorcategoria: \Tipo{\R}, factorPromo: \Tipo{\R}}
	{\R}
	{distancia * Factorcategoria * factorPromo}

	\vspace{1cm}
	\phantomsection\hypertarget{sec:EsNuevaTransaccion}{}\label{sec:EsNuevaTransaccion}
	\predLargo{EsNuevaTransaccion}{
		AM,AM_0 : \Tipo{AirMiles},idUsuario: \Tipo{\Z}, idReferidor: \Tipo{\Z}, \\TipoDeTransaccion: \Tipo{\Z}, Millas: \Tipo{\R}, distancia: \Tipo{\Z}, NombrePromo: \Tipo{\str}, FactorPromo: \Tipo{\R}
	}
	{AM.Historial = AM_0.Historial ++ \lbrack \tupla{idUsuario, idReferidor, \\ TipoDeTransaccion, Millas,NombrePromo, FactorPromo} \rbrack}
	{}

	\vspace{1cm}
	\phantomsection\hypertarget{sec:DarFactorCategoria}{}\label{sec:DarFactorCategoria}
	\auxLargo{DarFactorCategoria}
	{Historial: \Tipo{\seq{Transacciones}}, idUsuario: \Tipo{\Z}}
	{\Tipo{\R}}
	{
		\IfThenElse{
			\hyperref[sec:MillasTotales]{MillasTotales}(AM, idUsuario) < 10000
		}{
			1.0
			\\
		}{
			\IfThenElse{
				\hyperref[sec:MillasTotales]{MillasTotales}(AM, idUsuario) < 50000
			}{
				1.5
				\\
			}{
				\IfThenElse{
					\hyperref[sec:MillasTotales]{MillasTotales}(AM, idUsuario) < 1000000
				}{
					2.0
				}{
					0.0
				}
			}
		}
	}

	\vspace{0.5cm}
	\phantomsection\hypertarget{sec:DarCategoria}{}\label{sec:DarCategoria}
	\auxLargo{DarCategoria}
	{AM: \Tipo{Airmiles}, idUsuario: \Tipo{\Z}}
	{\Tipo{Categoria}}
	{
		\IfThenElse{
			\hyperref[sec:MillasTotales]{MillasTotales}(AM, idUsuario) < 10000
		}{
			Bronce
			\\
		}{
			\IfThenElse{
				\hyperref[sec:MillasTotales]{MillasTotales}(AM, idUsuario) < 50000
			}{
				Plata
				\\
			}{
				\IfThenElse{
					\hyperref[sec:MillasTotales]{MillasTotales}(AM, idUsuario) < 1000000
				}{
					Oro
				}{
					Platino
				}
			}
		}
	}


	\vspace{0.5cm}
	\noindent \comentario{En caso de que el usuario haga parte del viaje con promo y otra parte sin promo esta funcion busca esa diferencia}
	\\
	\phantomsection\hypertarget{sec:MillasACompletar}{}\label{sec:MillasACompletar}
	\\
	\auxLargo{MillasACompletar}
	{Promociones: \Tipo{\Diccionario{\str}{struct}},
		\\nombrePromo: \Tipo{\str}, distancia: \Tipo{\Z}, FactorCategoria: \Tipo{\R}}
	{\R}
	{
		\hyperref[sec:MillasCalculadas]{MillasCalculadas}(distancia,factorCategoria,\\
		Promociones[nombrePromo].factorDePromocion)-
		\\
		Promociones[nombrePromo].millasRestantes
	}

    \vspace{0.5cm}
    \phantomsection\hypertarget{sec:SumaVuelos}{}\label{sec:SumaVuelos}
    \auxLargo{SumaVuelos}
    {Vuelos: \seq{Vuelo}, Usuario: \Tipo{\Z}}
    {\Tipo{\Z}}
    {\sum\limits_{i=0}^{\norm{Vuelos}-1} 
    \IfThenElse{Usuario = Vuelos[i].Usuario}{Vuelos[i].Millas}{0}}

    \vspace{0.5cm}
    \noindent \comentario{Dado un usuario verifico las transferencias que recibio} \\
    \phantomsection\hypertarget{sec:SumaTransferenciasRecibidas}{}\label{sec:SumaTransferenciasRecibidas}
    \auxLargo{SumaTransferenciasRecibidas}
    {Transferencias: \seq{Transferencia}, IdUsuario: \Tipo{\Z}}
    {\Tipo{\Z}}
    {\sum\limits_{i=0}^{\norm{Transferencias}-1} 
    \IfThenElse{IdUsuario = Transferencias.UsuarioDestino \\}
    {Transferencias[i].Millas}
    {0}}

    \vspace{0.5cm}
    \noindent \comentario{Dado un usuario verifico las transferencias que envió} \\
    \phantomsection\hypertarget{sec:SumaTransferenciasEnviadas}{}\label{sec:SumaTransferenciasEnviadas}
    \auxLargo{SumaTransferenciasEnviadas}
    {Transferencias: \seq{Transferencia}, IdUsuario: \Tipo{\Z}}
    {\Tipo{\Z}}
    {\sum\limits_{i=0}^{\norm{Transferencias}-1} 
    \IfThenElse{IdUsuario = Transferencias.UsuarioOrigen \\}
    {Transferencias[i].Millas}
    {0}}

    \vspace{0.5cm}
    \noindent \comentario{Dado un usuario verifico las transferencias que recibió} \\
    \phantomsection\hypertarget{sec:SumaCanjeos}{}\label{sec:SumaCanjeos}
    \auxLargo{SumaCanjeos}
    {Canjeos: \seq{Canjeo}, IdUsuario: \Tipo{\Z}}
    {\Tipo{\Z}}
    {\sum\limits_{i=0}^{\norm{Canjeos}-1} 
    \IfThenElse{IdUsuario = Canjeos.Usuario}
    {Canjeos[i].Millas}
    {0}}

    \vspace{0.5cm}
    \noindent \comentario{Verifica que el usuario pertenezca a alguno de los campos de referidos y suma las millas} \\
    \phantomsection\hypertarget{sec:SumaReferidos}{}\label{sec:SumaReferidos}
    \auxLargo{SumaReferidos}
    {Referidos: \seq{Referido}, IdUsuario: \Tipo{\Z}}
    {\Tipo{\Z}}
    {\sum\limits_{i=0}^{\norm{Referidos}-1} 
    \IfThenElse{IdUsuario = Referidos.UsuarioReferidor \lor IdUsuario = Referidos.UsuarioReferido}
    {Referidos[i].Millas}
    {0}}

    \vspace{0.5cm}
    \phantomsection\hypertarget{sec:MillasTotales}{}\label{sec:MillasTotales}
    \noindent \comentario{Aca se suman todas las transacciones que recibe el usuario, y tambien a todos los que recibio y tambien si fue referido} \\
	\auxLargo{MillasTotales}
	{AM : \Tipo{Airmiles}, IdUsuario : \Tipo{\Z}}
	{\Tipo{\Z}}
	{res = SumaVuelos(AM.Vuelos,IdUsuario) + \\ 
           SumaTransferenciasRecibidas(AM.Transferencias, IdUsuario) + \\ 
           SumaReferidos(AM.Referidos,IdUsuario)}

    \vspace{0.5cm}
    \phantomsection\hypertarget{sec:MillasDisponibles}{}\label{sec:MillasDisponibles}
    \noindent \comentario{Se suma todo lo disponibles y tambien se restan los canjeos y transferencias que le hizo a otros} \\
	\auxLargo{MillasDisponibles}
	{AM : \Tipo{Airmiles}, IdUsuario : \Tipo{\Z}}
	{\Tipo{\Z}}
	{res = SumaVuelos(AM.Vuelos,IdUsuario) + \\ 
           SumaTransferenciasRecibidas(AM.Transferencias,IdUsuario) + \\
           SumaReferidos(AM.Referidos,IdUsuario) - \\
           SumaCanjeos(AM.Canjeos,IdUsuario) - \\
           SumaTransferenciasEnviadas(AM.Transferencias, IdUsuario)
    }

	\vspace{1.0cm}
	\phantomsection\hypertarget{sec:EsTransaccionSumadora}{}\label{sec:EsTransaccionSumadora}
	\predLargo{EsTransaccionSumadora}{Trans: Transacciones,idUsuario : \Tipo{\Z}}
	{ (Trans.TipoDeTransaccion = 0) \lor \\(IdUsuario = Trans.UsuarioDestino \land Trans.TipoDeTransaccion = 1)}


	\vspace{1.0cm}
	\phantomsection\hypertarget{sec:EsTransaccionDelUsuario}{}\label{sec:EsTransaccionDelUsuario}
	\pred{EsTransaccionDelUsuario}{Trans: Transacciones,idUsuario : \Tipo{\Z}}{\\Trans.UsuarioOrigen = idUsuario \lor Trans.UsuarioDestino = idUsuario}


	\vspace{0.5cm}
	\phantomsection\hypertarget{sec:SeConsumioLaPromocion}{}\label{sec:SeConsumioLaPromocion}
	\predLargo
	{SeConsumioLaPromocion}
	{
		nombrePromo: \Tipo{\str},
		AM: Airmiles
		AM_0: Airmiles
	}
	{AM.Promociones = DelKey(AM_0.Promociones,nombrePromo)}

	\vspace{0.5cm}

	\phantomsection\hypertarget{sec:MillasObtenidasPorVuelo}{}\label{sec:MillasObtenidasPorVuelo}
	\auxLargo{MillasObtenidasPorVuelo}
	{Historial: \Tipo{\seq{Transacciones}}, idOrigen: \Tipo{\Z}}
	{\Tipo{\R}}
	{
		\sum\limits_{i=0}^{\norm{Historial} - 1}
		\IfThenElse{
			\hyperref[sec:EsVuelo]{EsVuelo}(Historial[i],
			idOrigen)
		}{
			Historial[i].Millas
		}{
			0
		}
	}

	\vspace{1cm}

	\vspace{1cm}
	\phantomsection\hypertarget{sec:MillasDescontadasPorCanjeo}{}\label{sec:MillasDescontadasPorCanjeo}
	\auxLargo{MillasDescontadasPorCanjeo}
	{Historial: \Tipo{\seq{Transacciones}}, idOrigen: \Tipo{\Z}}
	{\Tipo{\R}}
	{
		\sum\limits_{i=0}^{\norm{Historial} - 1}
		\IfThenElse{
			\hyperref[sec:EsCanjeo]{EsCanjeo}(Historial[i],
			idOrigen)
		}{
			Historial[i].Millas
		}{
			0
		}
	}
	\vspace{1cm}
	\phantomsection\hypertarget{sec:EsVuelo}{}\label{sec:EsVuelo}
	\predLargo{EsVuelo}
	{Transaccion: \Tipo{Struct}, idOrigen: \Tipo{\Z}}
	{
		Transaccion.UsuarioOrigen = idOrigen \land
		\\
		Transaccion.TipoDeTransaccion = 0 \land
		\\
		Transaccion.UsuarioDestino = -1
	}

	\vspace{1cm}


	\phantomsection\hypertarget{sec:EsCanjeo}{}\label{sec:EsCanjeo}
	\predLargo{EsCanjeo}
	{Transaccion: \Tipo{struct}, idOrigen: \Tipo{\Z}}
	{
		Transaccion.UsuarioOrigen = idOrigen \land
		\\
		Transaccion.TipoDeTransaccion = 2 \land
		\\
		Transaccion.UsuarioDestino = -1
	}

	\vspace{1cm}
	\noindent \comentario{Es similar a EsTransferencia pero aca no importa el destino, unicamente chequeo que exista el usuario de origen.}
	\\
	\phantomsection\hypertarget{sec:EsTransferenciaEnviada}{}\label{sec:EsTransferenciaEnviada}
	\\
	\predLargo{EsTransferenciaEnviada}
	{Transaccion: \Tipo{Transacciones}, idOrigen: \Tipo{\Z}, Usuarios: \Tipo{\dict{\Z}{Categoria}}}
	{
		Transaccion.UsuarioOrigen = idOrigen \land
		\\
		Transaccion.TipoDeTransaccion = 1 \land
		\\
		Transaccion.UsuarioDestino \in claves(Usuarios)
	}

    \vspace{1cm}
	\phantomsection\hypertarget{sec:MillasPerdidas}{}\label{sec:MillasPerdidas}
	\auxLargo{MillasPerdidas}
	{AM: \Tipo{Airmiles},IdUsuario: \Tipo{\Z}}
	{\Tipo{\R}}
	{
        res = SumaTransferenciasEnviadas(AM.Transferencias,IdUsuario) +
        \\ SumaCanjeos(AM.Canjeos,IdUsuario)
    }

	\vspace{1cm}
	\phantomsection\hypertarget{sec:NoHayRepetidos}{}\label{sec:NoHayRepetidos}
	\predLargo{NoHayRepetidos}{res: \Tipo{\seq{\tupla{\Z,\Z}}}}
	{
		\paraTodo{t}{\tupla{\Tipo{\Z},\Tipo{\Z}}}{t \in res \implicaLuego \hyperref[sec:cantApariciones]{cantApariciones}(t,res)=1}
	}

	\phantomsection\hypertarget{sec:cantApariciones}{}\label{sec:cantApariciones}
	\auxLargo{cantApariciones}{t:\Tipo{\tupla{\Tipo{\Z,\Z}}},res: \Tipo{\seq{\Tipo{\Z},\Tipo{\Z}}}}{\Tipo{\Z}}{
		\sum\limits_{0}^{\norm{res}} \IfThenElse{t \in res}{1}{0}
	}
	\vspace{1cm}
	\phantomsection\hypertarget{sec:NoHayPermutacion}{}\label{sec:NoHayPermutacion}
	\predLargo{NoHayPermutacion}{res}{
		\paraTodo{t}{\Tipo{\tupla{\Tipo{\Z},\Tipo{\Z}}}}{t \in res \implicaLuego \tupla{{t_1},t_{0}} \notin res}
	}

	\vspace{0.5cm}
	\phantomsection\hypertarget{sec:TieneReciprocoHistorial}{}\label{sec:TieneReciprocoHistorial}
	\predLargo{TieneReciprocoHistorial}{Historial: \Tipo{\seq{Transacciones}},res: \seq{\tupla{\Tipo{\Z}, \Tipo{\Z}}}}
	{\paraTodo{x}{\tupla{\Tipo{\Z},\Tipo{\Z}}}{x \in res \implica \hyperref[sec:HayReciproco]{HayReciproco}(Historial,x)}}

	\vspace{0.5cm}
	\phantomsection\hypertarget{sec:HayReciproco}{}\label{sec:HayReciproco}
	\predLargo{HayReciproco}{Historial: \Tipo{\seq{Transacciones}},t: \Tipo{\Tipo{\tupla{\Tipo{\Z},\Tipo{\Z}}}}}
	{
		\existe{i}{\Tipo{\Z}}{0 \le i < \norm{Historial} \yLuego
			Historial[i].TipoDeTransaccion = 1 \yLuego
			\\
			(Historial[i].usuarioOrigen = t_{1} \land Historial[i].usuarioDestino = t_{0})}
	}

	\vspace{0.5cm}
	\phantomsection\hypertarget{sec:SonParejaExclusiva}{}\label{sec:SonParejaExclusiva}
	\predLargo{SonParejaExclusiva}{Historial: \Tipo{\seq{Transacciones}},res: \seq{\tupla{\Tipo{\Z}, \Tipo{\Z}}}}
	{
		\paraTodo{t}{\tupla{\Tipo{\Z},\Tipo{\Z}}}{ t \in res \implicaLuego
			\neg (\hyperref[sec:HayOtroUsuario]{HayOtroUsuario}(t0,t1,Historial) \lor
			\\
			\hyperref[sec:HayOtroUsuario]{HayOtroUsuario}(t1,t0,Historial))
		}
	}

	\phantomsection\hypertarget{sec:EstanTodosLosPares}{}\label{sec:EstanTodosLosPares}
	\predLargo{EstanTodosLosPares}{Usuarios: \Diccionario{id}{categoria},res}
	{
		\paraTodo{t}{\Tipo{\tupla{\Tipo{\Z},\Tipo{\Z}}}}{t0 \in claves(Usuarios) \land t1 \in claves(Usuarios) \implicaLuego t \in res }
	}

	\vspace{0.5cm}
	\phantomsection\hypertarget{sec:HayOtroUsuario}{}\label{sec:HayOtroUsuario}
	\predLargo{HayOtroUsuario}{t0,t1,Historial}
	{
		\paraTodo{i}{\Tipo{\Z}}{ 0 \le i < \norm{Historial} \yLuego Historial[i].UsuarioOrigen = t0 \implicaLuego
			\\
			Historial[i].UsuarioDestino \neq t1}
	}


\end{tad}

\end{document}
